\documentclass{article}
\usepackage[utf8]{inputenc}
\usepackage{graphicx}
\usepackage{amsmath}
\usepackage{amssymb}
\usepackage{geometry}
\usepackage{hyperref}
\usepackage{xurl}
\usepackage{cite}
\usepackage{indentfirst}
\usepackage{yhmath}
\usepackage{float}
\usepackage{listings}
\usepackage[dvipsnames]{xcolor}
\usepackage{booktabs} 
\usepackage{array}
\usepackage{titlesec}
\usepackage{multirow, multicol}
\geometry{letterpaper, margin=1in}
\hypersetup{hidelinks}
\def\UrlBreaks{%
    \do\/%
    \do\a\do\b\do\c\do\d\do\e\do\f\do\g\do\h\do\i\do\j\do\k\do\l%
    \do\m\do\n\do\o\do\p\do\q\do\r\do\s\do\t\do\u\do\v\do\w\do\x\do\y\do\z%
    \do\A\do\B\do\C\do\D\do\E\do\F\do\G\do\H\do\I\do\J\do\K\do\L%
    \do\M\do\N\do\O\do\P\do\Q\do\R\do\S\do\T\do\U\do\V\do\W\do\X\do\Y\do\Z%
    \do0\do1\do2\do3\do4\do5\do6\do7\do8\do9\do=\do/\do.\do:%
    \do\*\do\-\do\~\do\'\do\"\do\-}
\urlstyle{same}

\title{\textbf{ECE3700J Introduction to Computer Organization \\
Lab 3 Report --- Pipelined Processor}}

\author{\large Group 12 \\ % Update with your names maybe
% \large Student 1 Name (ID) \\
% \large Student 2 Name (ID) \\
% \large Student 3 Name (ID)
}

\date{\today} % Or \date{November 15, 2025}

\lstset{
    language=Verilog,
    backgroundcolor=\color{gray!20},
    basicstyle=\ttfamily\small,
    breakatwhitespace=false,
    breaklines=true,
    captionpos=b,
    keepspaces=true,
    numbers=left,
    numbersep=5pt,
    showspaces=false,
    showstringspaces=false,
    showtabs=false,
    tabsize=2,
    frame=single,
    framerule=0pt,
    framesep=3mm,
    xleftmargin=3mm,
    xrightmargin=3mm
}

\begin{document}

\maketitle

\newpage
\tableofcontents

\newpage
\section{Introduction}

This lab transitions from the single-cycle processor of Lab 2 to a 32-bit, 5-stage pipelined RISC-V processor. The pipelined architecture allows for greater instruction throughput by overlapping the execution of multiple instructions. The design is implemented in synthesizable Verilog HDL.

The processor supports an expanded subset of the RV32I instruction set, including:
\begin{itemize}
    \item \textbf{Arithmetic/Logical (R-type):} \texttt{add}, \texttt{sub}, \texttt{and}, \texttt{or}, \texttt{sll}, \texttt{srl}, \texttt{sra}
    \item \textbf{Arithmetic/Logical (I-type):} \texttt{addi}, \texttt{andi}, \texttt{slli}, \texttt{srli}, \texttt{srai}
    \item \textbf{Memory Access:} \texttt{lw}, \texttt{sw}, \texttt{lb}, \texttt{lbu}, \texttt{sb}
    \item \textbf{Conditional Branch:} \texttt{beq}, \texttt{bne}, \texttt{bge}, \texttt{blt}
    \item \textbf{Unconditional Jump:} \texttt{jal}, \texttt{jalr}
\end{itemize}

The primary goal is to correctly model the five pipeline stages (IF, ID, EX, MEM, WB) and the data/control paths connecting them. This design assumes a hazard-free test program, so hardware hazard detection, forwarding units, and stall logic are not implemented. The design's correctness is verified by simulating specific instructions and synthesizing the top-level module in Xilinx Vivado.

\section{Processor Architecture}

The processor implements a classic 5-stage RISC-V pipeline: Instruction Fetch (IF), Instruction Decode (ID), Execute (EX), Memory (MEM), and Write Back (WB). Four pipeline registers are used to pass data and control signals between stages: \texttt{if\_id\_reg}, \texttt{id\_ex\_reg}, \texttt{ex\_mem\_reg}, and \texttt{mem\_wb\_reg}.

\subsection{IF Stage (Instruction Fetch)}
The \texttt{PC} register holds the address of the current instruction. On each clock cycle, the \texttt{inst\_mem} is read at the address \texttt{PC}. The fetched instruction and \texttt{PC+4} are stored in the \texttt{if\_id\_reg}. The \texttt{PC} register itself is updated from a 3-way multiplexer (\texttt{pc\_branch\_selector}) which selects:
\begin{enumerate}
    \item \texttt{PC + 4} (default, sequential execution)
    \item \texttt{PC + imm} (for taken branches or \texttt{jal})
    \item \texttt{rs1 + imm} (for \texttt{jalr})
\end{enumerate}

\subsection{ID Stage (Instruction Decode)}
The \texttt{control} unit decodes the instruction (from \texttt{if\_id\_reg}) to generate all control signals. The \texttt{reg\_file} reads the values from source registers \texttt{rs1} and \texttt{rs2} concurrently. The \texttt{imm\_gen} module produces the sign-extended immediate value. All control signals, register values, the immediate, and PC are passed to the \texttt{id\_ex\_reg}.

\subsection{EX Stage (Execute)}
This stage performs calculations. The \texttt{alu\_src\_mux} selects either the \texttt{rs2} value or the immediate value as the second operand for the \texttt{alu}. The \texttt{alu\_control} unit generates the specific 4-bit ALU operation code. The \texttt{alu} performs the arithmetic or logical operation. A separate adder calculates the branch/jump target address. The \texttt{compare} module evaluates the branch condition (e.g., \texttt{zero} flag for \texttt{beq}/\texttt{bne}, ALU result for \texttt{blt}/\texttt{bge}).

\subsection{MEM Stage (Memory Access)}
This stage handles all data memory operations, receiving data from the \texttt{ex\_mem\_reg}.
\begin{itemize}
    \item For loads (\texttt{lw}, \texttt{lb}, \texttt{lbu}), \texttt{MemRead} is asserted, and the \texttt{data\_mem} is read at the address provided by the ALU result.
    \item For stores (\texttt{sw}, \texttt{sb}), \texttt{MemWrite} is asserted, and the \texttt{rs2} data (passed down from the ID stage) is written to memory.
\end{itemize}
The ALU result and any data read from memory are passed to the \texttt{mem\_wb\_reg}.

\subsection{WB Stage (Write Back)}
The final stage writes results back to the \texttt{reg\_file}. A multiplexer selects the data to be written based on control signals from the \texttt{mem\_wb\_reg}:
\begin{itemize}
    \item ALU result (for R-type or I-type arithmetic)
    \item Data from memory (for loads)
    \item \texttt{PC + 4} (for \texttt{jal} or \texttt{jalr} return address)
\end{itemize}
The \texttt{RegWrite} signal enables the write to the destination register \texttt{rd}.

\section{Instruction Simulation Results}

This section provides brief explanations and simulation waveforms for the key instructions specified in the lab manual. The waveforms demonstrate the flow of data and control signals through the 5-stage pipeline, incl. certain instruction lines from the \texttt{Lab3\_testcase.s} program.

\subsection{\texttt{add} Instruction}
\textbf{Explanation:} This waveform captures the Write-Back (WB) stage of the instruction \texttt{add t1, t0, t0} (fetched from PC \texttt{0x0C}). The cursor is at \texttt{t=42ns}, 4 cycles after its fetch at \texttt{t=18ns}.
\begin{itemize}
    \item \textbf{ID (at t=24ns):} Reads \texttt{rs1} (\texttt{t0}) and \texttt{rs2} (\texttt{t0}), both containing \texttt{0x193}.
    \item \textbf{EX (at t=30ns):} The ALU computes \texttt{0x193 + 0x193 = 0x326}.
    \item \textbf{WB (at t=42ns):} The cursor shows the result being written: \texttt{wb\_reg\_write} is \texttt{1}, \texttt{wb\_rd} is \texttt{06} (register \texttt{t1}), and \texttt{wb\_write\_data} is the correct result, \texttt{0x326}. The console at this time confirms \texttt{t1 = 0x00000326}.
\end{itemize}

\begin{figure}[H]
    \centering
    \includegraphics[width=0.9\linewidth]{Figures/add_sim.jpeg}
    \caption{Simulation of \texttt{add t1, t0, t0} (PC \texttt{0x0C}) in its Write-Back stage.}
    \label{fig:add_sim}
\end{figure}

\subsection{\texttt{addi} Instruction}
\textbf{Explanation:} This waveform shows the WB stage of the first instruction, \texttt{addi t0, x0, 0x193} (PC \texttt{0x00}). The instruction is fetched at \texttt{t=0ns} and reaches the WB stage at the cursor, \texttt{t=24ns}.
\begin{itemize}
    \item \textbf{EX (at t=12ns):} The ALU computes \texttt{rs1} (value \texttt{0} from \texttt{x0}) + the immediate \texttt{0x193}. The \texttt{alu\_result} is \texttt{0x193}.
    \item \textbf{WB (at t=24ns):} The WB signals are active: \texttt{wb\_reg\_write} is \texttt{1}, \texttt{wb\_rd} is \texttt{05} (register \texttt{t0}), and \texttt{wb\_write\_data} is \texttt{0x193}. The console confirms the write: \texttt{t0 = 0x00000193}.
\end{itemize}

\begin{figure}[H]
    \centering
    \includegraphics[width=0.9\linewidth]{Figures/addi_sim.jpeg} 
    \caption{Simulation of \texttt{addi t0, x0, 0x193} (PC \texttt{0x00}) in its Write-Back stage.}
    \label{fig:addi_sim}
\end{figure}

\subsection{\texttt{lw} Instruction}
\textbf{Explanation:} This waveform shows the WB stage of the \texttt{lw t4, 0(sp)} instruction (PC \texttt{0x110}). The cursor is at \texttt{t=426ns}.
\begin{itemize}
    \item \textbf{EX (at t=414ns):} The ALU calculates the memory address (\texttt{sp + 0}).
    \item \textbf{MEM (at t=420ns):} The \texttt{ex\_mem\_mem\_read} signal (not pictured) is \texttt{1}. The \texttt{data\_mem} is read, and the value \texttt{0xffcd8000} (which was stored by the \texttt{sw} instruction) is passed to the MEM/WB register.
    \item \textbf{WB (at t=426ns):} The \texttt{wb\_mem\_to\_reg} signal (not pictured) is \texttt{1}, selecting the data from memory. \texttt{wb\_reg\_write} is \texttt{1}, \texttt{wb\_rd} is \texttt{1d} (hex, or \texttt{29} decimal for \texttt{t4}), and \texttt{wb\_write\_data} is \texttt{ffcd8000}, the value loaded from memory.
\end{itemize}

\begin{figure}[H]
    \centering
    \includegraphics[width=0.9\linewidth]{Figures/lw_sim.jpeg} 
    \caption{Simulation of \texttt{lw t4, 0(sp)} (PC \texttt{0x110}) in its Write-Back stage.}
    \label{fig:lw_sim}
\end{figure}

\subsection{\texttt{sw} Instruction}
\textbf{Explanation:} This waveform captures the MEM stage of the \texttt{sw t0, 0(sp)} instruction (PC \texttt{0x108}), which occurs at \texttt{t=408ns}.
\begin{itemize}
    \item \textbf{EX (at t=402ns):} The ALU calculates the memory address (\texttt{sp + 0}).
    \item \textbf{MEM (at t=408ns):} The cursor is placed here. The key signals are in the EX/MEM register: \texttt{ex\_mem\_mem\_write} is \texttt{1}, and \texttt{ex\_mem\_write\_data} holds the value of \texttt{t0} (\texttt{0xffcd8000}), which is being written to data memory at this cycle.
    \item \textbf{WB (at t=414ns):} \texttt{wb\_reg\_write} is \texttt{0}, as \texttt{sw} does not write to the register file.
\end{itemize}

\begin{figure}[H]
    \centering
    \includegraphics[width=0.9\linewidth]{Figures/sw_sim.jpeg} 
    \caption{Simulation of \texttt{sw t0, 0(sp)} (PC \texttt{0x108}) in its Memory stage.}
    \label{fig:sw_sim}
\end{figure}

\subsection{\texttt{beq} Instruction}
\textbf{Explanation:} This waveform shows a \textit{taken} branch for the instruction \texttt{beq t2, t0, right\_branch\_3} (PC \texttt{0x80}). The instruction is in its EX stage at \texttt{t=204ns}.
\begin{itemize}
    \item \textbf{EX (at t=204ns):} The ALU compares \texttt{t2} and \texttt{t0} (both hold \texttt{0xffcd8000}). Since they are equal, \texttt{branch\_meet} is asserted to \texttt{1}.
    \item \textbf{Branch Logic:} Because \texttt{branch\_meet} is \texttt{1}, the \texttt{final\_sel} signal becomes \texttt{1}, selecting the branch target address.
    \item \textbf{IF (at t=210ns):} The \texttt{pc} register is updated to \texttt{0x9C} (the address of \texttt{right\_branch\_3}) instead of the sequential \texttt{0x84}, confirming the branch was successful.
\end{itemize}

\begin{figure}[H]
    \centering
    \includegraphics[width=0.9\linewidth]{Figures/beq_sim.jpeg} 
    \caption{Simulation of a taken \texttt{beq} instruction (PC \texttt{0x80}) in its Execute stage.}
    \label{fig:beq_sim}
\end{figure}

\subsection{\texttt{jal} Instruction}
\textbf{Explanation:} This waveform captures the WB stage of the \texttt{jal x1, memory\_test} instruction (PC \texttt{0xF8}) at \texttt{t=396ns}. This instruction performs two actions: jumping and linking.
\begin{itemize}
    \item \textbf{The Jump (at t=384ns):} In the EX stage, the \texttt{jump} signal is \texttt{1} and \texttt{final\_sel} becomes \texttt{1}. At the next clock edge (\texttt{t=390ns}), the \texttt{pc} is updated to the jump target, \texttt{0x108} (the address of \texttt{memory\_test}).
    \item \textbf{The Link (at t=396ns):} The cursor shows the return address being saved. \texttt{wb\_reg\_write} is \texttt{1}, \texttt{wb\_rd} is \texttt{01} (register \texttt{x1}), and \texttt{wb\_write\_data} is \texttt{0xFC}, which is the correct return address (\texttt{PC+4}, or \texttt{0xF8 + 4}).
\end{itemize}

\begin{figure}[H]
    \centering
    \includegraphics[width=0.9\linewidth]{Figures/jal_sim.jpeg} 
    \caption{Simulation of \texttt{jal x1, memory\_test} (PC \texttt{0xF8}) in its Write-Back stage.}
    \label{fig:jal_sim}
\end{figure}

\newpage
\section{RTL Schematic}

The Verilog design was synthesized using Xilinx Vivado to verify its synthesizability and examine the hardware structure. The resulting RTL schematic (Figure~\ref{fig:rtl}) shows the high-level blocks corresponding to the pipeline stages and major components like the \texttt{reg\_file}, \texttt{alu}, and \texttt{data\_mem}, connected by the pipeline registers.

\begin{figure}[H]
    \centering
    \includegraphics[width=\linewidth]{Figures/RTL_schematic_pipelined.jpeg} 
    \caption{Vivado-generated RTL schematic of the 5-Stage Pipelined Processor.}
    \label{fig:rtl}
\end{figure}

\section{Source Code Listing}
\label{sec:source}

The complete Verilog source code (including \texttt{top} design and all instantiated submodules) is provided in the submission \texttt{lab3.zip} on Canvas.

\section{Conclusion}

A 32-bit, 5-stage pipelined RISC-V processor was successfully designed, implemented, and verified in Verilog. The processor supports a broader range of instructions than in the previous lab, including arithmetic, memory, branch, and jump operations.

Key outcomes:
\begin{itemize}
    \item A 5-stage pipeline (IF, ID, EX, MEM, WB) was implemented, with pipeline registers passing data and control signals between stages.
    \item The datapath and control logic were designed to handle R, I, S, B, and J-type instruction formats.
    \item Simulation of key instructions (\texttt{add}, \texttt{addi}, \texttt{lw}, \texttt{sw}, \texttt{beq}, \texttt{jal}) confirmed correct data flow, control signal generation, and pipeline operation for each instruction type.
    \item The design was successfully synthesized in Vivado, confirming the RTL is synthesizable and producing a schematic that matches the intended pipelined architecture.
\end{itemize}

This lab demonstrates the fundamental principles of pipelined processor design, showing how instruction execution can be overlapped to improve performance.

\newpage
\begin{thebibliography}{99} 
\bibitem{Lab3} ECE3700J Teaching Team. \texttt{T6 Pipelined Processor.pdf}, \texttt{Lab3.pdf}, \texttt{Lab3\_testcase.s}, and \texttt{Lab3\_testcase.txt}. Retrieved on 27 Oct 2025 from \hyperlink{site}{oc.sjtu.edu.cn}.
\end{thebibliography}

\end{document}